\documentclass[a4paper,10pt]{article}

% --- PAQUETES DE DISEÑO ---
\usepackage[utf8]{inputenc}
\usepackage[T1]{fontenc}
\usepackage[default]{lato} 
\usepackage[margin=1.25cm]{geometry} 
\usepackage{xcolor}
\usepackage{hyperref}
\usepackage{fontawesome5} 
\usepackage{titlesec}
\usepackage{enumitem}

% --- COLORES ---
\definecolor{primary}{RGB}{40, 100, 160} 
\definecolor{secondary}{RGB}{80, 80, 80} 

% --- ESTILO DE TÍTULOS ---
\titleformat{\section}
{\Large\bfseries\color{primary}}
{}{0em}
{}[\titlerule]
\titlespacing*{\section}{0pt}{1.2em}{0.6em}

% --- QUITAR NUMERACIÓN ---
\pagestyle{empty}

% --- COMIENZO DEL DOCUMENTO ---
\begin{document}

% === CABECERA ===
\begin{center}
    % APLICAMOS escape_tex AL NOMBRE Y TÍTULO
    {\Huge \textbf{Tu Nombre Completo}} \\ \vspace{0.2cm}
    {\Large \textcolor{secondary}{Perfil Multidisciplinar}} \\ \vspace{0.2cm}
    
    % Iconos de contacto
    % NOTA: En los enlaces (\href), el primer argumento es la URL (mejor no escapar salvo que de error),
    % el segundo argumento es el texto visible (SÍ hay que escapar).
    \small
            \faEnvelope \ \href{mailto:hola@tudominio.com}{hola@tudominio.com} \quad | \quad
                \faLinkedin \ \href{https://linkedin.com/in/tuusuario}{LinkedIn} \quad | \quad
                \faGithub \ \href{https://github.com/tuusuario}{GitHub} \quad | \quad
                \faGlobe \ \href{https://tu-portfolio.github.io}{Portfolio}
    \end{center}

\vspace{0.3cm}

% === PERFIL ===
\section*{Perfil}
% Aquí aplicamos el filtro al texto del perfil
Profesional híbrido con experiencia en desarrollo backend e investigación en psicología. Apasionado por la intersección entre datos, conducta humana y tecnología.

% === EXPERIENCIA ===
\section*{Experiencia Laboral}
\begin{itemize}[leftmargin=0cm, label={}]
        \item 
        % Aplicamos filtro a puesto, fecha, empresa, ubicación y descripción
        \textbf{\large Backend Developer} \hfill \textbf{2021 - Presente} \\
        \textcolor{primary}{\textit{Tech Solutions S.L.}} \ \textbullet \ \textit{Madrid, España} \\
        \vspace{0.05cm}
        Desarrollo de microservicios con Python y FastAPI. Optimización de consultas SQL.
        \vspace{0.3cm}
        \item 
        % Aplicamos filtro a puesto, fecha, empresa, ubicación y descripción
        \textbf{\large Soporte Técnico} \hfill \textbf{2016 - 2018} \\
        \textcolor{primary}{\textit{HelpDesk Corp}} \ \textbullet \ \textit{Remoto} \\
        \vspace{0.05cm}
        Resolución de incidencias de nivel 2 y trato directo con clientes internacionales.
        \vspace{0.3cm}
    \end{itemize}

% === FORMACIÓN ===
\section*{Formación Académica}
\begin{itemize}[leftmargin=0cm, label={}]
        \item 
        \textbf{Grado en Psicología} \hfill 2014 - 2018 \\
        \textit{Universidad X}
        \vspace{0.15cm}
        \item 
        \textbf{Bootcamp Full Stack Web} \hfill 2020 \\
        \textit{Code School Y}
        \vspace{0.15cm}
    \end{itemize}

% === HABILIDADES ===
\section*{Habilidades y Herramientas}
\begin{itemize}[leftmargin=0cm, label={}]
    \item 
            % Filtro también aquí por si pones "C++" o "C#"
        \colorbox{primary!10}{\textcolor{black}{\textbf{Python / Django}}} \quad
            % Filtro también aquí por si pones "C++" o "C#"
        \colorbox{primary!10}{\textcolor{black}{\textbf{Git - GitHub}}} \quad
    \end{itemize}

\end{document}